\documentclass[11pt]{article}
\usepackage{amsmath,amssymb,amsthm}
\usepackage[margin=2.5cm]{geometry}

\newtheorem{theorem}{Theorem}
\newtheorem{lemma}[theorem]{Lemma}
\newtheorem{proposition}[theorem]{Proposition}
\newtheorem{corollary}[theorem]{Corollary}
\theoremstyle{remark}
\newtheorem{remark}[theorem]{Remark}
\theoremstyle{definition}
\newtheorem*{statement}{Official statement}

\newcommand{\pr}{\operatorname{pr}}

\title{Problem 4, Cell 3: Sun's conjecture for all $k\le 8$}
\author{}
\date{}

\begin{document}
\maketitle

\begin{statement}
For integers $a$ and $m\ge1$ write $a\ (\mathrm{mod}\ m)=\{a+mt:t\in\mathbb{Z}\}$. A finite family
of classes $a_1\ (\mathrm{mod}\ m_1),\dots,a_k\ (\mathrm{mod}\ m_k)$ is pairwise disjoint if
$a_i\ (\mathrm{mod}\ m_i)\cap a_j\ (\mathrm{mod}\ m_j)=\emptyset$ whenever $i<j$. Nothing is assumed
about the moduli beyond $m_i\ge1$: they may repeat, and the residues $a_i$ may repeat as well. By the
Chinese remainder theorem,
\[
  a_i\ (\mathrm{mod}\ m_i)\cap a_j\ (\mathrm{mod}\ m_j)\neq\emptyset\iff\gcd(m_i,m_j)\mid a_i-a_j.\tag{$*$}
\]
\emph{Question.} Let $k\ge2$ and let $a_1\ (\mathrm{mod}\ m_1),\dots,a_k\ (\mathrm{mod}\ m_k)$ be
pairwise disjoint. Must there be a pair $i<j$ with $\gcd(m_i,m_j)\ge k$?

A first range of sizes: after C1 and C2, the sizes $k=5,6,7$ and $8$ remain. Prove the statement
for every $k\le8$.
\end{statement}

\noindent Theorem~\ref{thm:main} below is exactly this statement, for all $2\le k\le8$ (the question assumes
$k\ge2$). The notation $a\bmod m$ below is the class $a\ (\mathrm{mod}\ m)$ of the statement.
Lemma~\ref{lem:crt} is $(*)$, proved for completeness.

\paragraph{Known result.} The question is Sun's disjoint congruence classes conjecture~\cite{Sun06}.
O'Bryant~\cite{OBr07} proved it for all $k\le20$, so this case is known. The proof below is an
independent proof and is self-contained.

For an integer $m\ge1$ and $a\in\mathbb{Z}$ write $a\bmod m=\{x\in\mathbb{Z}:x\equiv a\pmod m\}$.
A family of residue classes is \emph{disjoint} if its members are pairwise disjoint.

\begin{theorem}\label{thm:main}
Let $2\le k\le 8$ and let $a_1\bmod m_1,\dots,a_k\bmod m_k$ be a disjoint family of residue
classes. Then $\gcd(m_i,m_j)\ge k$ for some $i\ne j$.
\end{theorem}

The proof is by hand throughout; no computation is used. The cases $k=5,6$ follow from
Lemma~\ref{lem:pp}, and the cases $k=7,8$ from Proposition~\ref{prop:78}. The same arguments
also give $k=2,3,4$ again.

\medskip\noindent\textbf{Notation.} Throughout, $a_1\bmod m_1,\dots,a_k\bmod m_k$ ($k\ge2$) is a
disjoint family, $[k]=\{1,\dots,k\}$, and $g_{ij}=\gcd(m_i,m_j)$ for $i\ne j$. For an integer
$n\ge 1$, $\pr(n)$ denotes the set of primes dividing $n$.

\section{Preliminaries}

\begin{lemma}\label{lem:crt}
Let $m,n\ge1$, $a,b\in\mathbb{Z}$ and $g=\gcd(m,n)$. Then $(a\bmod m)\cap(b\bmod n)\ne\emptyset$ if
and only if $g\mid a-b$.
\end{lemma}

\begin{proof}
($\Rightarrow$) If $x\equiv a\pmod m$ and $x\equiv b\pmod n$, then $a\equiv x\equiv b\pmod g$.
($\Leftarrow$) Choose $u,v\in\mathbb{Z}$ with $um+vn=g$ (B\'ezout). Write $a-b=gt$ and put
$x=a-umt$. Then $x\equiv a\pmod m$, and $x=a-(g-vn)t=b+vnt\equiv b\pmod n$.
\end{proof}

\begin{corollary}\label{cor:basic}
For all $i\ne j$ we have $g_{ij}\ge2$ and $a_i\not\equiv a_j\pmod{g_{ij}}$. Consequently, if
$g_{ij}\mid q$ for some integer $q$, then $a_i\not\equiv a_j\pmod q$.
\end{corollary}

\begin{proof}
By Lemma~\ref{lem:crt}, disjointness means $g_{ij}\nmid a_i-a_j$, and this forces $g_{ij}\ne1$.
If $a_i\equiv a_j\pmod q$ and $g_{ij}\mid q$, then $a_i\equiv a_j\pmod{g_{ij}}$, a contradiction.
\end{proof}

We use the pigeonhole principle in the following form. If integers $a_i$ ($i\in I$) are
pairwise incongruent modulo $q$, then $|I|\le q$.

\paragraph{Prime channels.}
For $i\in[k]$ let
\[
  \pi(i)=\{\,p \text{ prime}: p\mid g_{ij}\text{ for some } j\ne i\,\},
  \qquad S_p=\{\,i\in[k]: p\in\pi(i)\,\}.
\]

\begin{lemma}[Channel lemma]\label{lem:channel}
\begin{enumerate}
\item[(a)] $\pr(g_{ij})=\pi(i)\cap\pi(j)$ for all $i\ne j$; in particular $\pi(i)\cap\pi(j)\ne\emptyset$.
\item[(b)] $|S_p|\ne1$ for every prime $p$.
\item[(c)] If $p$ is a prime and $i\ne j$ are both in $S_p$, then $p\mid g_{ij}$.
\end{enumerate}
\end{lemma}

\begin{proof}
(a) The inclusion $\subseteq$ holds by definition of $\pi$. Conversely, if $p\in\pi(i)\cap\pi(j)$,
then $p$ divides some $g_{il}$ and hence $m_i$, and similarly $p\mid m_j$; so $p\mid g_{ij}$. The set
$\pr(g_{ij})$ is nonempty because $g_{ij}\ge2$ (Corollary~\ref{cor:basic}).
(b) If $i\in S_p$, then $p\mid g_{ij}$ for some $j\ne i$, and then $j\in S_p$ as well.
(c) This is (a).
\end{proof}

\section{All gcds prime powers: $k\le 6$}

\begin{lemma}[Prime-power lemma]\label{lem:pp}
Let $2\le M\le5$ and suppose every $g_{ij}$ is a prime power with $g_{ij}\le M$. For each prime
$p\le M$ let $q_p$ be the largest power of $p$ with $q_p\le M$, and let $P$ be the set of primes
$\le M$. Then
\[
  k\le\max\Bigl(\max_{p\in P}q_p,\ \tbinom{|P|}{2}\Bigr).
\]
\end{lemma}

\begin{proof}
By Lemma~\ref{lem:channel}(a), each $\pi(i)\cap\pi(j)=\pr(g_{ij})$ is a single prime, which lies in $P$.
Every prime in $\pi(i)$ divides some $g_{ij}\le M$, so $\pi(i)\subseteq P$.

\emph{Case I: some prime $p$ lies in every $\pi(i)$.} Then $p\in\pi(i)\cap\pi(j)=\pr(g_{ij})$ for all $i\ne j$,
so every $g_{ij}$ is a power of $p$ that is at most $M$. Hence $g_{ij}\mid q_p$. By
Corollary~\ref{cor:basic}, the $a_i$ are pairwise incongruent modulo $q_p$, so $k\le q_p$ by pigeonhole.

\emph{Case II: no prime lies in every $\pi(i)$.} If some $\pi(i)=\{p\}$, then every $\pi(j)$ ($j\ne i$) meets
$\{p\}$, so $p$ lies in all $\pi(\cdot)$, which is Case I. Hence $|\pi(i)|\ge2$ for all $i$. Two distinct
indices $i\ne j$ cannot satisfy $\pi(i)=\pi(j)$, since then $|\pi(i)\cap\pi(j)|\ge2$. Also
$|\pi(i)|\le2$ for every $i$. Indeed, $|P|\le3$ because $M\le5$. So $|\pi(i)|\ge3$ would force
$\pi(i)=P\supseteq\pi(j)$ for any $j\ne i$ (such a $j$ exists since $k\ge2$), and then
$|\pi(i)\cap\pi(j)|=|\pi(j)|\ge2$. Hence the $\pi(i)$ are pairwise distinct $2$-element subsets
of $P$, so $k\le\binom{|P|}{2}$.
\end{proof}

\begin{corollary}\label{cor:k56}
Theorem~\ref{thm:main} holds for $2\le k\le6$.
\end{corollary}

\begin{proof}
For $k=2$ this is Corollary~\ref{cor:basic}. Let $3\le k\le6$ and suppose all $g_{ij}\le k-1$. Put
$M=k-1\in\{2,3,4,5\}$. Every integer in $\{2,\dots,5\}$ is a prime power, so Lemma~\ref{lem:pp}
applies:
\begin{center}
\begin{tabular}{c|c|c|c|c}
$k$ & $M$ & $P$ & $\max_p q_p$ & $\binom{|P|}{2}$\\ \hline
3 & 2 & $\{2\}$ & 2 & 0\\
4 & 3 & $\{2,3\}$ & 3 & 1\\
5 & 4 & $\{2,3\}$ & 4 & 1\\
6 & 5 & $\{2,3,5\}$ & 5 & 3
\end{tabular}
\end{center}
In every row the bound of Lemma~\ref{lem:pp} equals $k-1<k$, a contradiction.
\end{proof}

\begin{remark}
For $k=5$ the sets $S_2$ and $S_3$ play the role of the sets $E$ and $T$ of Cell~2 (in general
only $S_2\subseteq E$ and $S_3\subseteq T$). Case II would require
two indices with $\pi(i)=\pi(j)=\{2,3\}$, i.e.\ $6\mid g_{ij}$. In Case I with $p=2$ all gcds
divide $4$, and with $p=3$ all gcds equal $3$.
\end{remark}

\section{Gcds in $\{2,3,4,6\}$}

\begin{lemma}\label{lem:W}
If every $g_{ij}$ lies in $\{2,3,4,6\}$, then $k\le6$.
\end{lemma}

\begin{proof}
Let $E=\{i:2\mid m_i\}$ and $T=\{i:3\mid m_i\}$. Each $i$ lies in $E\cup T$, because $g_{ij}$ (for any
$j\ne i$) is divisible by $2$ or $3$ and divides $m_i$. If $i\in E\setminus T$ and $j\in T\setminus E$, then
$g_{ij}$ is coprime to $6$, which is impossible. Hence $T\setminus E=\emptyset$ or $E\setminus T=\emptyset$, that is,
$[k]=E$ or $[k]=T$. Put $A=E\cap T=\{i:6\mid m_i\}$. For $i\ne j$ in $A$ we have $6\mid g_{ij}$, so
$g_{ij}=6$.

\emph{Case $[k]=T$.} Let $B=T\setminus E$. If $i\ne j$ and at least one of them is in $B$, then $3\mid g_{ij}$
and $2\nmid g_{ij}$, so $g_{ij}=3$. By Corollary~\ref{cor:basic}:
\begin{itemize}
\item the $a_b$ ($b\in B$) are pairwise incongruent mod $3$;
\item each $a_i$ ($i\in A$) is incongruent mod $3$ to every $a_b$ ($b\in B$);
\item the $a_i$ ($i\in A$) are pairwise incongruent mod $6$.
\end{itemize}
So the residues $a_i\bmod6$ ($i\in A$) are distinct, and they lie in the preimage of the
$3-|B|$ classes of $\mathbb{Z}/3$ not used by $B$. Each class of $\mathbb{Z}/3$ has exactly two
preimages in $\mathbb{Z}/6$, so $|A|\le2(3-|B|)$. Then $k=|A|+|B|\le6-|B|\le6$.

\emph{Case $[k]=E$.} Let $B=E\setminus T$, $A_2=\{i\in A:4\mid m_i\}$ and $A_1=A\setminus A_2$. We have
$|A_2|\le1$, since two indices of $A_2$ would give $12\mid g_{ij}$. For $b\ne b'$ in $B$, $g_{bb'}$ is
even and not divisible by $3$, so $g_{bb'}\in\{2,4\}$ and $g_{bb'}\mid4$. For $i\in A_1$ and $b\in B$,
$g_{ib}$ is even, not divisible by $3$ (as $3\nmid m_b$) and not divisible by $4$ (as $4\nmid m_i$), so
$g_{ib}=2$. By Corollary~\ref{cor:basic}:
\begin{itemize}
\item the $a_b$ ($b\in B$) are pairwise incongruent mod $4$;
\item each $a_i$ ($i\in A_1$) has parity opposite to every $a_b$ ($b\in B$);
\item the $a_i$ ($i\in A$) are pairwise incongruent mod $6$.
\end{itemize}
If $B=\emptyset$, then $[k]=A$ and $k\le6$ by pigeonhole mod $6$. If the $a_b$ ($b\in B$) have both
parities, then $A_1=\emptyset$ and $k\le|B|+|A_2|\le4+1=5$. Otherwise all $a_b$ have the same parity
$\varepsilon$. Then $|B|\le2$, since there are only two classes mod $4$ of parity $\varepsilon$. Every
$a_i$ with $i\in A_1$ has parity $1-\varepsilon$, and these $a_i$ are pairwise incongruent mod $6$. Only
three classes mod $6$ have a given parity, so $|A_1|\le3$. Altogether $k\le2+3+1=6$.
\end{proof}

\section{Gcds at most $7$: the cases $k=7,8$}

\begin{proposition}\label{prop:78}
Suppose every $g_{ij}\le7$. Then $k\le7$. If moreover no $g_{ij}$ equals $7$, then $k\le6$.
\end{proposition}

\begin{proof}
For $n\in\{2,\dots,7\}$ the sets $\pr(n)$ are $\{2\},\{3\},\{2\},\{5\},\{2,3\},\{7\}$. By
definition of $\pi(i)$, since every $g_{ij}\le7$, every $\pi(i)\subseteq\{2,3,5,7\}$. Let $p\in\{5,7\}$. For $i\ne j$ in $S_p$,
Lemma~\ref{lem:channel}(c) gives $p\mid g_{ij}\le7$, so
\begin{equation}\label{eq:Sp}
  g_{ij}=p\quad\text{and}\quad\pi(i)\cap\pi(j)=\{p\}\qquad(i\ne j\in S_p).
\end{equation}
In particular the sets $\pi(i)\setminus\{p\}$, $i\in S_p$, are pairwise disjoint.

\emph{Step 1. For $p\in\{5,7\}$: either $S_p=[k]$, in which case $k\le p$ and every $g_{ij}=p$; or
$|S_p|\le|\pi(l)|$ for every $l\notin S_p$.}
If $S_p=[k]$, then \eqref{eq:Sp} gives $g_{ij}=p$ for all pairs, and $k\le p$ by
Corollary~\ref{cor:basic} and pigeonhole. Otherwise, let $l\notin S_p$. For each $i\in S_p$, the set
$\pi(i)\cap\pi(l)$ is nonempty (Lemma~\ref{lem:channel}(a)) and contained in $\pi(i)\setminus\{p\}$
(since $p\notin\pi(l)$). These nonempty sets are pairwise disjoint subsets of $\pi(l)$, so
$|S_p|\le|\pi(l)|$.

Let $O=[k]\setminus(S_5\cup S_7)$. Thus $\pi(l)\subseteq\{2,3\}$ for $l\in O$.

\emph{Step 2. The case $O=\emptyset$.} Then $[k]=S_5\cup S_7$.
\begin{itemize}
\item If $S_5=[k]$, then all $g_{ij}=5$ and $k\le5$.
\item If $S_7=[k]$, then all $g_{ij}=7$ and $k\le7$.
\item Otherwise, for each $p\in\{5,7\}$ some $l\notin S_p$ exists. Since $\pi(l)\subseteq\{2,3,5,7\}\setminus\{p\}$,
  Step~1 gives $|S_p|\le3$. Hence $k\le|S_5|+|S_7|\le6$.
\end{itemize}

\emph{Step 3. The case $O\ne\emptyset$.} Fix $l_0\in O$ and let $p\in\{5,7\}$ with $S_p\ne\emptyset$.
By Step~1 (note $l_0\notin S_p$), $|S_p|\le|\pi(l_0)|\le2$, and by Lemma~\ref{lem:channel}(b),
$|S_p|=2$. Write $S_p=\{i,i'\}$. For every $l\in O$, the sets $\pi(i)\cap\pi(l)$ and $\pi(i')\cap\pi(l)$
are nonempty, disjoint (by \eqref{eq:Sp}) and contained in $\pi(l)\subseteq\{2,3\}$. Therefore
$\pi(l)=\{2,3\}$. After renaming $i,i'$ we may take $\pi(i)\cap\pi(l)=\{2\}$ and $\pi(i')\cap\pi(l)=\{3\}$.
The names $i,i'$ do not depend on $l$: we have $\pi(i)\cap\{2,3\}=\pi(i)\cap\pi(l)$ for every
$l\in O$.

\emph{Step 3(a): $S_5=S_7=\emptyset$.} Then $O=[k]$, and every $\pr(g_{ij})=\pi(i)\cap\pi(j)\subseteq\{2,3\}$.
With $g_{ij}\le7$ this means $g_{ij}\in\{2,3,4,6\}$. Lemma~\ref{lem:W} gives $k\le6$.

\emph{Step 3(b): some $S_p=\{i,i'\}$ is nonempty, with notation as above.} We derive
constraints on the $a_l$ ($l\in O$).
\begin{itemize}
\item For $l\ne l'$ in $O$: $\pr(g_{ll'})=\{2,3\}$ and $g_{ll'}\le7$, so $g_{ll'}=6$. The $a_l$ are
  pairwise incongruent mod $6$.
\item For $l\in O$: $\pr(g_{i'l})=\{3\}$, so $g_{i'l}=3$ and $a_l\not\equiv a_{i'}\pmod3$.
\item For $l\in O$: $\pr(g_{il})=\{2\}$, so $g_{il}\in\{2,4\}$. At most one $l\in O$ has $4\mid m_l$, since
  two such $l,l'$ would give $12\mid g_{ll'}$. For every other $l\in O$ we get $g_{il}=2$ (as
  $g_{il}\mid m_l$), hence $a_l\not\equiv a_i\pmod2$.
\end{itemize}
Let $O'=\{l\in O:4\nmid m_l\}$, so $|O|\le|O'|+1$. The residues $a_l\bmod6$ ($l\in O'$) are distinct.
Each has parity different from $a_i$ and residue mod $3$ different from $a_{i'}$. By the Chinese
remainder theorem, $\mathbb{Z}/6\cong\mathbb{Z}/2\times\mathbb{Z}/3$, and exactly $1\cdot2=2$ classes
mod $6$ meet both conditions. So $|O'|\le2$ and $|O|\le3$. Step~3 also showed $|S_q|\in\{0,2\}$ for
both $q\in\{5,7\}$. Hence
\[
  k\le|S_5|+|S_7|+|O|\le2+2+3=7.
\]

\emph{Conclusion.} In every case $k\le7$. Now suppose no $g_{ij}$ equals $7$. Then no prime channel is
$7$, i.e.\ $S_7=\emptyset$ by \eqref{eq:Sp} and Lemma~\ref{lem:channel}(b). The cases give:
\begin{itemize}
\item Step~2: $k\le5$ ($S_5=[k]$ is forced, because $S_7=\emptyset$ and $O=\emptyset$);
\item Step~3(a): $k\le6$;
\item Step~3(b): $k\le|S_5|+|O|\le2+3=5$.
\end{itemize}
So $k\le6$.
\end{proof}

\section{Proof of Theorem~\ref{thm:main}}

The cases $2\le k\le6$ are Corollary~\ref{cor:k56}. Suppose $k=7$ and all $g_{ij}\le6$. Then
Proposition~\ref{prop:78} (second part) gives $k\le6$, a contradiction. Suppose $k=8$ and all
$g_{ij}\le7$. Then Proposition~\ref{prop:78} (first part) gives $k\le7$, a contradiction. \qed

\begin{remark}[Sharpness]
For every $k$ the family $r\bmod k$ ($0\le r<k$) is disjoint with all pairwise gcds equal to
$k$. The bounds in Lemma~\ref{lem:W} and Proposition~\ref{prop:78} are attained. The six classes
$r\bmod6$ have all gcds equal to $6$. The seven classes $r\bmod7$ have all gcds equal to $7$.
\end{remark}

\begin{remark}[Computer sanity check, not part of the proof]
The script \texttt{check\_p4\_c3.py} checks the theorem numerically; the proof does not use it.
\begin{itemize}
\item \emph{Graph.} For each $k\in\{5,6,7,8\}$ it forms all classes $a\bmod m$ with $0\le a<m\le60$
  (1829 classes). Two classes are joined when they are disjoint and their moduli have gcd $<k$.
\item \emph{Disjointness.} This is checked literally, from the definition, and not via $(*)$ or
  Lemma~\ref{lem:crt}. The classes $a\bmod m$ and $b\bmod n$ meet iff $b-a\equiv mt\pmod n$ for some
  $t$, and the set $\{mt\bmod n:0\le t<n\}$ is listed explicitly.
\item \emph{Result.} An exact branch-and-bound search finds a maximum clique size of exactly $k-1$
  in each case (5.5\,s in total). This agrees with the theorem.
\end{itemize}
\end{remark}

\begin{thebibliography}{9}
\bibitem{OBr07} K.~O'Bryant, \emph{On Z.-W.\ Sun's disjoint congruence classes conjecture}, in
  \emph{Combinatorial Number Theory}, de Gruyter, Berlin, 2007, pp.~403--411; arXiv:math/0604347.
\bibitem{Sun06} Z.-W.~Sun, \emph{Finite covers of groups by cosets or subgroups}, Internat.\ J.\ Math.\
  \textbf{17} (2006), 1047--1064; arXiv:math/0501451.
\end{thebibliography}

\end{document}
